% --- IMPORTS ---
\documentclass[leqno]{article}
\usepackage{graphicx}
\usepackage{dsfont}
\usepackage{mathtools}
\usepackage{amssymb}
\usepackage{amsmath}
\usepackage{amsthm}
\usepackage{float}
\usepackage{ragged2e}
\usepackage{tensor}
\usepackage{fancyhdr}
\usepackage{empheq}
\usepackage[most]{tcolorbox}
\usepackage{enumitem}
\usepackage{dirtytalk}
\usepackage{marginnote}
\usepackage{microtype}
\usepackage[
  a4paper,
  left=0.8in,
  right=1in,
  top=1in,
  bottom=0.5in,
  headheight=14pt,
  headsep=0.4in
]{geometry}
\setlength{\parindent}{0cm}
% --- TikZ Packages ---
\usepackage{pgfplots}
\pgfplotsset{compat=1.18}
\usepgfplotslibrary{fillbetween, polar}
\usetikzlibrary{calc, positioning}
\usetikzlibrary{angles, quotes}
\usetikzlibrary{arrows.meta}
\usetikzlibrary{decorations.pathreplacing, decorations.markings}
\usetikzlibrary{patterns}
\usetikzlibrary{intersections}
% --- URL Packages ---
\usepackage[colorlinks=true,linkcolor=red,citecolor=magenta,urlcolor=black]{hyperref}%
\urlstyle{same}
% --- END OF IMPORTS ---

% --- CUSTOM COMMANDS ---
\RenewDocumentCommand{\marks}{o m}{%
  \IfValueTF{#1}%
    {\marginnote{\raggedright\textbf{[#2]}}[#1]}%
    {\marginnote{\raggedright\textbf{[#2]}}}%
}
\newcommand{\vspacerule}[2]{%
  \par\vspace*{#1}%
  \noindent\hrulefill\par
  \vspace*{#2}%
}
\definecolor{scarlet}{RGB}{205, 40, 75}
\newcommand{\questionitem}{%
  \item\phantomsection
  \addcontentsline{toc}{section}{Question \number\value{enumi}}%
}
\renewcommand{\contentsname}{Questions}
% --- END OF CUSTOM COMMANDS ---

% --- HEADER CONFIG ---
\pagestyle{fancy}
\fancyhead{}
\fancyfoot{}
\renewcommand{\headrulewidth}{0.9pt}
\lhead{Invariant Points and Lines}
\chead{}
\rhead{\href{https://danieljsmith.org}{danieljsmith.org}}
% --- END OF HEADER CONFIG ---

\begin{document}
\vspace*{20pt}
\tableofcontents
\newpage
\begin{enumerate}[
  label=\textbf{\arabic*.},
  leftmargin=*,
  topsep=0pt,
  partopsep=0pt,
  parsep=0pt
]


% ---- Question 1 ----
\questionitem

The matrix $M$ represents the transformation $T$ and is given by
\[
M = \begin{pmatrix}
p + 3 & -2 \\[6pt]
 p + 2 & 1 - p
\end{pmatrix}
\] \\
\begin{enumerate}[label=\textbf{(\alph*)}]
  \item In the case when $p = -1$ show that the image of the point $(3, 4)$ under $T$ is the point $(-2, 11)$. \marks{2} \\
  \item In the case when $p = 2$ find the gradient of the line of invariant points under $T$. \marks{3} \\
\end{enumerate}

\newpage

% ---- Question 2 ----
\questionitem

A linear transformation of the plane is represented by
\[
\begin{pmatrix}
1 & 4 \\
2 & 3
\end{pmatrix}
\]
Determine the invariant line with negative gradient. \marks{4}

\newpage

% ---- Question 3 ----
\questionitem

The matrix $A$ is given by
\[
A = \frac{1}{5}\begin{pmatrix}-3 & -4 \\ -4 & 3\end{pmatrix}
\] \\
You are given that $A$ represents the transformation $T$, where $T$ is a reflection in a straight line through the origin, $O$. \\
\begin{enumerate}[label=\textbf{(\alph*)}]
  \item Explain why $T$ must have a line of invariant points. State the geometric significance of this line. \marks{2} \\
  \item By considering the line of invariant points for $T$, determine the equation of the mirror line.

Give your answer in the form $y = mx + c$. \marks{4} \\
  \item The coordinates of the point $P$ are $(4, 7)$. \\
  By considering the image of $P$ under the transformation $T$, or otherwise, determine the coordinates of the point on the mirror line which is closest to $P$. \marks{3} \\
  \item The line with equation $2x + by = 6$ is an invariant line for $T$. \\
  Determine the value of $b$. \marks{2} \\
\end{enumerate}

\newpage

% ---- Question 4 ----
\questionitem

The matrix $M$ is given by
\[
M = \begin{pmatrix} 1 & 0 \\[6pt] 0 & 6 \end{pmatrix}\begin{pmatrix} \frac{1}{2} & \frac{\sqrt{3}}{2} \\[6pt] \frac{\sqrt{3}}{2} & -\frac{1}{2} \end{pmatrix}
\] \\
\begin{enumerate}[label=\textbf{(\alph*)}]
  \item The matrix $M$ represents a sequence of two geometrical transformations in the $x$-$y$ plane. \\
  Give full details of the two transformations, stating clearly the order in which they are applied. \marks{4} \\
  \item Write $M^{-1}$ as the product of two matrices, neither of which is $I$. \marks{2} \\
  \item Find the equations of the invariant lines, through the origin, of the transformation represented by $M$. \marks{5} \\
\end{enumerate}

\newpage


% ---- Question 5 ----
\questionitem

A linear transformation of the plane is represented by the matrix
\[
\mathbf{M}=\begin{pmatrix}2 & -1\\ p & 0\end{pmatrix}
\]
where $p$ is a constant. \\
\begin{enumerate}[label=\textbf{(\alph*)}]
  \item Show that the line $x=0$ is not invariant. For a line through the origin with equation $y=mx$, show that it is invariant only if
  \[
  m^2-2m+p=0
  \]
  Hence find the set of values of $p$ for which the transformation has no invariant lines through the origin. \marks{5}\\
  \item Given that the transformation multiplies areas by $3$ and reverses orientation, find the invariant lines. \marks{3} \\
\end{enumerate}

\newpage

% ---- Question 6 ----
\questionitem

The $2 \times 2$ matrix $\mathbf{A}$ represents a transformation $\mathbf{T}$ which has the following properties. \\

\begin{itemize}
\item The image of the point $(1,-1)$ is the point $(2,0)$.
\item An object shape whose area is $4$ is transformed to an image shape whose area is $24$.
\item $\mathbf{T}$ has a line of invariant points. \\
\end{itemize}

\begin{enumerate}[label=(\roman*)]
  \item Find a possible matrix for $\mathbf{A}$. \marks{8} \\
\end{enumerate}

The transformation $\mathbf{S}$ is represented by the matrix $\mathbf{B}$ where
\[
\mathbf{B} = \begin{pmatrix}2 & 1 \\ 2 & 3\end{pmatrix}
\] \\
\begin{enumerate}[label=(\roman*), resume]
  \item Find the equation of the line of invariant points of $\mathbf{S}$. \marks{2} \\
  \item Show that any line of the form $y = 2x + c$ is an invariant line of $\mathbf{S}$. \marks{3} \\
\end{enumerate}

\newpage

% ---- Question 7 ----
\questionitem

The matrix
\[
B = \begin{pmatrix} 1 & -4 \\ 2 & -1 \end{pmatrix}
\]
represents the linear transformation $T$. \\
Any invariant line for $T$ must be parallel to an invariant line through the origin. \\
\begin{enumerate}[label=\textbf{(\alph*)}]
  \item A line through the origin has equation $y=mx$, where $m$ is real. Show that, for $m \ne \tfrac{1}{4}$, its image under $T$ has equation
  \[
  y = \frac{2-m}{1-4m}x
  \]
  State also the image of the line $y=\tfrac{1}{4}x$ and the image of the vertical line $x=0$. \marks{3}\\
  \item Hence prove that $T$ has no invariant lines. \marks{2} \\
\end{enumerate}

\newpage

% ---- Question 8 ----
\questionitem

The matrix of a linear transformation $T$ is
\[
M = \begin{pmatrix} p+2 & -5 \\[6pt] -2 & p-1 \end{pmatrix}
\] \\
\begin{enumerate}[label=\textbf{(\alph*)}]
  \item Find the values of $p$ for which $T$ is not invertible. \marks{2} \\
\end{enumerate}
It is now given that $p = 3$, so that
\[
M = \begin{pmatrix} 5 & -5 \\ -2 & 2 \end{pmatrix}
\] \\
\begin{enumerate}[label=\textbf{(\alph*)}, resume]
  \item Find the equations of the invariant lines, through the origin, of the transformation represented by $M$. \marks{6} \\
\end{enumerate}

\newpage

% ---- Question 9 ----
\questionitem

\[
A = \begin{pmatrix} b & a - 3 \\[6pt] 2a - 1 & -2 \end{pmatrix}
\]
where $a$ and $b$ are non-zero constants. \\
Given that the matrix $A$ is self-inverse, \\
\begin{enumerate}[label=\textbf{(\alph*)}]
  \item determine the value of $b$ and the possible values of $a$. \marks{5} \\
\end{enumerate}
The matrix $A$ represents a linear transformation $M$ of the plane. \\
Using the smaller value of $a$ from part (a), \\
\begin{enumerate}[label=\textbf{(\alph*)}, resume]
  \item show that the invariant points of the linear transformation $M$ form a line, stating the equation of this line. \marks{3} \\
\end{enumerate}
\end{enumerate}
\end{document}
